\section{Conclusion}


We have presented the 4-Loop Self-Improvement Architecture, a structured
telemetry-driven system for autonomous coding agent self-improvement. The
architecture's central contributions are:

\begin{itemize}[leftmargin=1.5em]
    \item The telemetry thesis: immutable structured records of tool invocations
          are a superior substrate for self-improvement compared to narrative
          summaries, with a formal $\Omega(\log n)$ entropy advantage.

    \item The Build-It-Now protocol: deterministic triggers at $\theta=3$
          achieve F1 = 0.91 compared to LLM self-judgment F1 = 0.66, while
          eliminating backlog accumulation.

    \item Four-loop structure with clean separation of concerns: each loop
          operates at a well-defined timescale, reads from and writes to
          well-defined stores, and never pollutes the active context window.

    \item Convergence guarantees: under mild conditions, the 4-Loop system
          converges to unit tool success rate almost surely.
\end{itemize}

The architecture is practically grounded in the bmo results (11 tools,
100 sessions) and theoretically grounded in information theory and
stochastic convergence. It is available as a TypeScript plugin for the
Hanzo bot SDK, with full telemetry service, friction detector, active
learning hooks, and maintenance tooling.

We believe that the move from aspirational to telemetric self-improvement
is a necessary step for deploying autonomous coding agents that reliably
improve with use. The failure modes of narrative-based approaches
(backlog growth, context pollution, LLM hallucination, preference drift)
are not edge cases but fundamental properties of the paradigm. Structured
telemetry and deterministic triggers address these failure modes at their
root.

% ============================================================================
% REFERENCES
% ============================================================================
\bibliographystyle{plainnat}
\begin{thebibliography}{99}

\bibitem[bmo(2024)]{bmo2024}
ngrok Engineering.
\newblock bmo: A self-improving coding agent.
\newblock Technical report, ngrok Inc., 2024.
\newblock URL \url{https://ngrok.com/blog/bmo}.

\bibitem[Brown et al.(2020)]{brown2020language}
T.~Brown, B.~Mann, N.~Ryder, M.~Subbiah, J.~D. Kaplan, P.~Dhariwal,
  A.~Neelakantan, P.~Shyam, G.~Sastry, A.~Askell, et~al.
\newblock Language models are few-shot learners.
\newblock In \emph{Advances in Neural Information Processing Systems},
  volume~33, pages 1877--1901. Curran Associates, Inc., 2020.

\bibitem[Chen et al.(2021)]{chen2021evaluating}
M.~Chen, J.~Tworek, H.~Jun, Q.~Yuan, H.~P. de~Oliveira~Pinto, J.~Kaplan,
  H.~Edwards, Y.~Burda, N.~Joseph, G.~Brockman, et~al.
\newblock Evaluating large language models trained on code.
\newblock \emph{arXiv preprint arXiv:2107.03374}, 2021.

\bibitem[Cover and Thomas(2006)]{cover2006elements}
T.~M. Cover and J.~A. Thomas.
\newblock \emph{Elements of Information Theory}.
\newblock Wiley-Interscience, 2nd edition, 2006.

\bibitem[Guo et al.(2017)]{guo2017calibration}
C.~Guo, G.~Pleiss, Y.~Sun, and K.~Q. Weinberger.
\newblock On calibration of modern neural networks.
\newblock In \emph{Proceedings of the 34th International Conference on Machine
  Learning}, pages 1321--1330. PMLR, 2017.

\bibitem[Hanzo AI(2026a)]{hanzo2026aso}
Hanzo AI Research.
\newblock Active semantic optimization: Training-free adaptation via Bayesian
  product-of-experts decoding.
\newblock Technical report, Hanzo AI Inc., February 2026.

\bibitem[Hanzo AI(2026b)]{hanzo2026dso}
Hanzo AI Research.
\newblock Decentralized semantic optimization: Byzantine-robust prior
  aggregation for collective model adaptation.
\newblock Technical report, Hanzo AI Inc., February 2026.

\bibitem[Hinton(2002)]{hinton2002training}
G.~E. Hinton.
\newblock Training products of experts by minimizing contrastive divergence.
\newblock \emph{Neural Computation}, 14\penalty0 (8):\penalty0 1771--1800,
  2002.

\bibitem[Jimenez et al.(2024)]{jimenez2024swebench}
C.~E. Jimenez, J.~Yang, A.~Wettig, S.~Yao, K.~Pei, O.~Press, and K.~Narasimhan.
\newblock SWE-bench: Can language models resolve real-world GitHub issues?
\newblock In \emph{Proceedings of the 12th International Conference on Learning
  Representations}, 2024.

\bibitem[Kadavath et al.(2022)]{kadavath2022language}
S.~Kadavath, T.~Conerly, A.~Askell, T.~Henighan, D.~Drain, E.~Perez,
  N.~Schiefer, Z.~Hatfield-Dodds, N.~DasSarma, E.~Tran-Johnson, et~al.
\newblock Language models (mostly) know what they know.
\newblock \emph{arXiv preprint arXiv:2207.05221}, 2022.

\bibitem[Levenshtein(1966)]{levenshtein1966binary}
V.~I. Levenshtein.
\newblock Binary codes capable of correcting deletions, insertions, and
  reversals.
\newblock \emph{Soviet Physics Doklady}, 10\penalty0 (8):\penalty0 707--710,
  1966.

\bibitem[Madaan et al.(2023)]{madaan2023selfrefine}
A.~Madaan, N.~Tandon, P.~Gupta, S.~Hallinan, L.~Gao, S.~Wiegreffe,
  U.~Alon, N.~Dziri, S.~Prabhumoye, Y.~Yang, et~al.
\newblock Self-refine: Iterative refinement with self-feedback.
\newblock In \emph{Advances in Neural Information Processing Systems},
  volume~36, 2023.

\bibitem[Qin et al.(2024)]{qin2024toolllm}
Y.~Qin, S.~Liang, Y.~Ye, K.~Zhu, L.~Yan, Y.~Lu, Y.~Lin, X.~Cong, X.~Tang,
  B.~Qian, et~al.
\newblock ToolLLM: Facilitating large language models to master 16000+
  real-world APIs.
\newblock In \emph{Proceedings of the 12th International Conference on Learning
  Representations}, 2024.

\bibitem[Schick et al.(2024)]{schick2024toolformer}
T.~Schick, J.~Dwivedi-Yu, R.~Dess{\`{\i}}, R.~Raileanu, M.~Lomeli,
  L.~Zettlemoyer, N.~Cancedda, and T.~Scialom.
\newblock Toolformer: Language models can teach themselves to use tools.
\newblock In \emph{Advances in Neural Information Processing Systems},
  volume~36, 2024.

\bibitem[Shannon(1948)]{shannon1948mathematical}
C.~E. Shannon.
\newblock A mathematical theory of communication.
\newblock \emph{The Bell System Technical Journal}, 27\penalty0 (3):\penalty0
  379--423, 1948.

\bibitem[Shinn et al.(2023)]{shinn2023reflexion}
N.~Shinn, F.~Cassano, E.~Berman, A.~Gopinath, K.~Narasimhan, and S.~Yao.
\newblock Reflexion: Language agents with verbal reinforcement learning.
\newblock In \emph{Advances in Neural Information Processing Systems},
  volume~36, 2023.

\bibitem[Sutton and Barto(2018)]{sutton2018reinforcement}
R.~S. Sutton and A.~G. Barto.
\newblock \emph{Reinforcement Learning: An Introduction}.
\newblock MIT Press, 2nd edition, 2018.

\bibitem[Wang et al.(2023)]{wang2023voyager}
G.~Wang, Y.~Xie, Y.~Jiang, A.~Mandlekar, C.~Xiao, Y.~Zhu, L.~Fan, and
  A.~Anandkumar.
\newblock Voyager: An open-ended embodied agent with large language models.
\newblock \emph{arXiv preprint arXiv:2305.16291}, 2023.

\bibitem[Wang et al.(2024)]{wang2024openhands}
X.~Wang, B.~Li, Y.~Song, F.~F. Xu, X.~Tang, M.~Zhuge, J.~Pan, Y.~Song,
  B.~Li, J.~Singh, et~al.
\newblock Openhands: An open platform for AI software developers as generalist
  agents.
\newblock \emph{arXiv preprint arXiv:2407.16741}, 2024.

\bibitem[Wu et al.(2023)]{wu2023autogen}
Q.~Wu, G.~Bansal, J.~Zhang, Y.~Wu, B.~Li, E.~Zhu, L.~Jiang, X.~Zhang,
  S.~Zhang, J.~Liu, et~al.
\newblock AutoGen: Enabling next-gen LLM applications via multi-agent
  conversation.
\newblock \emph{arXiv preprint arXiv:2308.08155}, 2023.

\end{thebibliography}

